% Speaker Notes for CS265 Midterm Presentation
% Compile: xelatex slides.tex
\documentclass[11pt]{article}
\usepackage[margin=1in]{geometry}
\usepackage{booktabs}
\usepackage{enumitem}
\usepackage{hyperref}
\usepackage{xcolor}
\definecolor{question}{RGB}{0,100,180}
\definecolor{note}{RGB}{100,100,100}
\newcommand{\slide}[1]{\subsection*{Slide: #1}}
\newcommand{\say}[1]{\textbf{Say:} #1}
\newcommand{\qa}[2]{\textcolor{question}{\textbf{Q:} #1} \\ \textbf{A:} #2}
\newcommand{\timing}[1]{\hfill\textcolor{note}{\textit{(#1)}}}

\title{Speaker Notes\\{\large CS265 Activation Checkpointing --- Midterm Presentation}}
\author{Hadi Khalaf}
\date{}

\begin{document}
\maketitle

\textbf{Total time:} 25--30 minutes. Pace: $\sim$1 minute per slide. The results slides and key findings can be expanded if time permits.

\tableofcontents

\section{Motivation (Slides 1--6)} \timing{$\sim$6 min}

\slide{Title Slide}
\say{This project implements activation checkpointing for PyTorch --- a technique that reduces GPU memory by trading off extra computation. I'll walk through what we built and what we found.}

\slide{Outline}
\say{Three sections: motivation, Phase 1 (the profiler), and Phase 2 (the selection algorithm). I'll finish with results and next steps.}

\slide{The Memory Problem in DNN Training}
\say{When you train a neural network, the GPU has to hold four things simultaneously: the model parameters (weights), the activations (intermediate outputs from each layer), the gradients, and the optimizer state like Adam's moment buffers. The key point is that activations typically account for 70--85\% of peak memory. And the problem is getting worse --- model sizes grow exponentially but GPU memory grows linearly.}

\qa{Why 70--85\%? Where does that number come from?}{It's empirical --- from analyzing memory breakdowns of standard models like ResNet, BERT, GPT. Parameters and optimizer state are proportional to model size, but activations are proportional to model size \emph{times} batch size \emph{times} sequence length. For large batches, activations dominate.}

\qa{What about techniques like mixed precision or gradient accumulation?}{Those help too, but they're orthogonal. Mixed precision halves the bytes per element; AC reduces the number of elements stored simultaneously. They compose well --- you can use both.}

\slide{Why Activations Dominate Peak Memory}
\say{Here's the key insight. The forward pass produces activations Z1 through Z4, left to right. But the backward pass needs them in \emph{reverse} order. So Z1 is produced first but needed last --- it sits in memory the longest. At the peak (typically the end of the forward pass), ALL activations are alive simultaneously. That's why they dominate.}

\qa{Can't you just free each activation as soon as the next layer is done with it?}{No --- the backward pass needs them later. For example, to compute the gradient of a linear layer, you need its forward input. PyTorch keeps these alive in autograd closures. They can't be freed until the corresponding backward hook fires.}

\slide{Activation Checkpointing: The Idea}
\say{The idea is simple: don't keep all activations in memory. Discard some after the forward pass, and recompute them during backward when needed. This trades compute for memory. The question is: which ones to discard? You want to discard the ones that are cheap to recompute (like ReLU) and keep the expensive ones (like convolutions). That's what Phase 2 decides.}

\qa{Isn't recomputation wasteful? You're doing the same work twice.}{Yes, it's extra compute. But on modern GPUs, the forward pass is fast --- typically a few milliseconds per operation. The memory savings (hundreds of MBs) can let you increase batch size, which improves training throughput more than the recomputation costs.}

\slide{Project Overview}
\say{The project has three phases. Phase 1 builds the profiler that collects the data. Phase 2 implements the mu-TWO algorithm that makes the selection decisions. Phase 3, which is next, will actually rewrite the graph to insert recomputation nodes. I've completed Phases 1 and 2, tested on four models.}


\section{Phase 1: Graph Profiler (Slides 7--17)} \timing{$\sim$10 min}

\slide{Background: Why FX Tracing?}
\say{My first milestone used PyTorch module hooks to instrument the model from the outside. That worked for profiling, but hooks can't rewrite the graph. For Phase 3, we need to insert new operations into the computation. PyTorch's FX framework gives us a single, symbolic graph that contains the entire training step --- forward, backward, and optimizer --- and it's rewritable. So we switched to FX.}

\qa{What's the difference between FX tracing and torch.compile?}{torch.compile uses Dynamo for tracing, which handles more Python constructs. FX's \texttt{make\_fx} is simpler --- it uses fake tensors to trace through the ops symbolically. The course starter code uses \texttt{make\_fx}, so that's what we build on.}

\slide{The Traced Graph Structure}
\say{The traced graph has a clear structure. Placeholder nodes at the top are the inputs: model parameters, optimizer states, and batch data. Then the forward pass operations --- transpose, matmul, relu, etc. Then a separator node, then loss computation, then a second separator, then backward operations, and finally the optimizer step. The key point: this is a \emph{single flat graph} containing the entire training iteration. The two separator nodes are what let us partition it into four regions.}

\say{Each block here is a group of FX nodes. Every individual operation --- each relu, each matrix multiply --- is its own node with explicit inputs and users. That's what makes lifetimes directly readable.}

\qa{What are the separator nodes? Why do we need them?}{Without them, the graph is one continuous stream of ATen operations --- \texttt{sum}, \texttt{mm}, \texttt{copy\_} all look the same. There's no reliable way to tell where forward ends and backward begins. \texttt{SEPFunction} is a custom \texttt{torch.autograd.Function} whose forward calls \texttt{separator.sep(x)} (returns x unchanged) and whose backward calls \texttt{separator.sep\_backward(grad)} (returns grad unchanged). When \texttt{make\_fx} traces through \texttt{loss.backward()}, it records both the forward and backward calls as nodes in the graph. So you get: \texttt{... relu\_9, sum\_1, sep, view, ones\_like, sep\_backward, threshold\_backward ...}. The two marker nodes are no-ops mathematically but are explicit, searchable boundaries in the graph.}

\qa{Could you just search for the first gradient operation instead?}{Not reliably. Some operations like \texttt{view} or \texttt{ones\_like} appear in the loss region between forward and backward, and they look identical to forward ops. The separator gives a guaranteed boundary regardless of loss function or model architecture.}

\slide{Phase 1 Architecture}
\say{The profiler has three components. Static analysis happens once when created --- walks the graph for boundaries, tensor types, and lifetimes. Runtime profiling measures timing and memory per node during actual execution. The visualizer produces the stacked bar chart.}

\slide{Static Analysis: Boundary Detection}
\say{Simple scan for separator targets. Everything before sep is forward, after sep\_backward is backward. The first \_foreach operation after backward marks the optimizer region.}

\slide{Static Analysis: Parameter Identification}
\say{This was one of the trickiest challenges. All three types of placeholder --- parameters, optimizer states, and batch data --- look identical in the graph. They're all \texttt{placeholder} nodes, and FakeTensors all have \texttt{requires\_grad=False}, so we can't use that flag. And with \texttt{foreach=True} Adam, the optimizer is decomposed into hundreds of small \texttt{\_foreach\_*} ops --- there's no single \texttt{\_fused\_adam} node whose arguments list the parameters.}

\say{The insight is to look at \emph{where} each placeholder is used across the graph regions. From the debug output: a weight matrix like \texttt{arg0\_1} has users in the forward pass (\texttt{t} at step 19) \emph{and} in the optimizer (\texttt{\_foreach\_addcdiv} at step 171). The batch tensor \texttt{arg3\_3} has users in forward (\texttt{addmm}) and backward (\texttt{mm\_2}), but \emph{not} the optimizer --- it's never updated. An optimizer state like \texttt{arg2\_1} has users \emph{only} in the optimizer region.}

\say{So the rule is: parameters are the only placeholders that appear in both forward and optimizer. Two lines of code: check for a forward user and an optimizer user. This works for every model --- DummyModel, ResNet18, ResNet50, BERT --- regardless of optimizer type.}

\qa{What if a model uses SGD instead of Adam?}{The usage-pattern heuristic still works --- parameters are always consumed in forward (to compute outputs) and updated in the optimizer step (gradient descent). The specific decomposition doesn't matter.}

\qa{What about shared parameters (e.g., tied embeddings in a transformer)?}{Shared parameters would still satisfy both conditions --- they'd have forward users (wherever they're used) and optimizer users (the update step). They'd appear once as a placeholder and be correctly identified.}

\qa{Why not just check which placeholders have \texttt{requires\_grad=True}?}{We tried that first. It doesn't work because \texttt{make\_fx} traces with \texttt{FakeTensorMode}, which creates fake tensors that all have \texttt{requires\_grad=False} regardless of the original tensor's setting. The gradient flag isn't propagated through symbolic tracing.}

\slide{Five Tensor Categories}
\say{The course spec requires five categories. Let me walk through each.}

\say{PARAM --- we just discussed this. Placeholder nodes with users in both the forward and optimizer regions. The model weights and biases.}

\say{ACT --- the intermediate activations that AC targets. These are \texttt{call\_function} nodes produced during the forward pass that have at least one user in the backward pass. For example, \texttt{relu} at step 85 is consumed by \texttt{threshold\_backward} at step 119. The ``forward producer, backward consumer'' pattern is what defines an activation --- it's a tensor sitting idle in memory waiting for backward.}

\say{GRAD --- gradient tensors. We can only identify these when \texttt{fused=True} Adam is used, because the \texttt{\_fused\_adam} node has structured arguments: \texttt{args[0]} = parameters, \texttt{args[1]} = gradients. With \texttt{foreach=True}, gradients are mixed into the decomposed \texttt{\_foreach\_*} ops and we can't distinguish them. That's why the output shows GRAD count = 0 --- they get lumped into OTHER. This doesn't affect AC decisions since we only need PARAM and ACT for that.}

\say{OPTIMIZER\_STATE --- Adam's \texttt{exp\_avg}, \texttt{exp\_avg\_sq}, and \texttt{step} counters. Placeholder nodes whose users are \emph{exclusively} in the optimizer region. Unlike parameters which are also used in forward, optimizer states are only touched during the Adam update.}

\say{OTHER --- everything else: batch data (forward + backward users but not optimizer), the loss tensor, backward computation nodes like \texttt{threshold\_backward} and \texttt{mm}, and all the optimizer computation nodes like \texttt{\_foreach\_add} and \texttt{copy\_}.}

\qa{Does it matter that GRAD count is 0?}{Not for our purposes. The AC algorithm only needs to know which nodes are activations (to consider evicting) and which are parameters (to know they're always available for recomputation). Gradient classification is informational for the memory breakdown chart but doesn't affect decisions.}

\slide{Intermediate Activations and Lifetimes}
\say{Four conditions define an intermediate activation. For each one, we compute the idle period: the gap between last forward use and first backward use. Longer idle period + larger size = better checkpointing candidate.}

\qa{Are all forward-pass tensors intermediates?}{No. Only those with backward users. A tensor used only in forward (e.g., the output of a pooling layer that feeds the loss but not backward directly) wouldn't qualify. In practice most forward computation nodes do feed into backward.}

\slide{Runtime Profiling}
\say{The static analysis reads the graph structure. Runtime profiling actually \emph{executes} it. The FX Interpreter calls \texttt{run\_node()} for every node --- we override it to measure two things.}

\say{First, timing. We record a CUDA event before and after the node executes. \texttt{torch.cuda.synchronize()} is critical here --- GPU operations are asynchronous, meaning Python returns immediately while the GPU kernel is still running. Without sync, we'd measure Python dispatch time, not actual GPU kernel time. \texttt{elapsed\_time} between the two events gives us the real wall-clock GPU time.}

\say{Second, memory. We read \texttt{torch.cuda.memory\_allocated()} before and after. The delta tells us how much GPU memory the operation allocated or freed. For example, \texttt{addmm} (a Linear forward) allocates a new output tensor (+400 KB), while some backward ops free tensors they no longer need.}

\say{The activation swapping is the clever part. After the static analysis identifies each intermediate's last forward use and first backward use, we build a swap schedule. When we execute the node at \texttt{last\_fwd\_access} for intermediate X, we move X to CPU with \texttt{.cpu()} right after. When we reach \texttt{first\_bwd\_access}, we move it back with \texttt{.cuda()} right before. This does two things: it measures actual PCIe swap latency per tensor, and it prevents OOM --- without swapping, profiling ResNet50 or BERT would keep all intermediates in GPU memory simultaneously. The swap happens \emph{outside} the \texttt{super().run\_node()} call, so it doesn't contaminate the per-node timing.}

\qa{Doesn't synchronizing after every node slow down profiling?}{Yes, significantly. Synchronization forces the GPU to finish all pending work before Python continues. In normal training, operations overlap and pipeline. During profiling we sacrifice throughput for measurement accuracy. That's why we only profile for a few iterations (3 measurement runs), not the entire training.}

\qa{Why swap to CPU instead of just deleting the tensor?}{We need the tensor back for the backward pass --- the interpreter environment (\texttt{self.env}) must have it when a backward node tries to consume it. Swapping preserves the data while freeing GPU memory. Deleting would cause a crash when backward tries to access it.}

\slide{Memory Breakdown Visualization}
\say{The stacked bar chart. Green is activations --- they accumulate during forward, get freed during backward. The peak tells you where to focus. If green dominates, AC helps. If red (optimizer state) dominates, you need a different strategy.}

\slide{Phase 1 Results: Profiler Accuracy}
\say{Validated against actual GPU measurements. ResNet18 is 0.97x --- near-perfect. DummyModel is 0.45x because CUDA allocator rounding is proportionally large on small tensors. BERT is 1.18x because we overcount optimizer temporaries. All within reasonable bounds.}

\qa{Why does the static estimate differ from actual GPU?}{Our estimate uses \texttt{numel * element\_size} from FakeTensor metadata. The GPU also has: CUDA allocator block rounding (512 bytes or 2 MB), kernel temporary buffers invisible to FX, autograd SavedVariable wrappers, and fragmentation.}


\section{Phase 2: AC Selection Algorithm (Slides 18--28)} \timing{$\sim$10 min}

\slide{The Subset Selection Problem}
\say{Classic optimization: N items, each with a cost and a benefit. Which to select? NP-hard in general. We use the greedy heuristic from mu-TWO.}

\slide{mu-TWO: Background}
\say{mu-TWO is an MLSys 2023 paper for multi-model training. Its full algorithm considers both swap and recompute. We use only recompute --- appropriate for single-model AC. The ranking metric is recompute ratio: bytes freed per millisecond of extra compute.}

\qa{Why not implement swap too?}{Swap requires modeling PCIe bandwidth and overlap with computation from a second model (the multiplexer). For single-model training without a second model to overlap with, swap offers less benefit. The profiler does measure swap times for future extension.}

\slide{The Greedy Algorithm (Algorithm B)}
\say{The algorithm: compute baseline peak, set target, check if AC can help at all, then greedily evict the highest-ratio candidate. After each eviction, re-simulate peak. Stop when target is met or no more progress. The memory simulator is the key --- it tells us exactly when we've done enough.}

\slide{Memory Simulator}
\say{The simulator walks the timeline. Retained tensors are live from production to last use. Evicted tensors are live during forward, FREE during idle, briefly live at recomputation. Peak is the max over all steps. This is why ``memory freed'' exceeds ``peak reduction'' --- evicted tensors are still live in forward.}

\slide{Cascading Recomputation}
\say{When A is evicted and B depends on A, recomputing B now requires first recomputing A. B's cost goes up, ratio goes down. This prevents runaway chains. From Algorithms E and F in the paper.}

\slide{Early Termination: The BERT Finding}
\say{BERT's peak is at the optimizer step. No activations alive there. Our quick check detects this and returns immediately. The lesson: not every model benefits from AC. BERT's bottleneck is the optimizer, not activations.}

\qa{How would you fix BERT's memory?}{\texttt{fused=True} Adam (single node instead of hundreds of foreach ops), optimizer state sharding (like ZeRO), or gradient accumulation to reduce batch-level optimizer state.}

\slide{Validation: Input Availability}
\say{Every recomputed tensor's inputs must be available. Iterate until stable --- moving one back to retained can enable another eviction.}

\slide{Results: ResNet18}
\say{Peak 506 to 431 MB. Only 8 evictions needed --- all ReLUs. All convolutions retained. 1 ms extra cost.}

\slide{Results: ResNet50 (bs=16)}
\say{This is the biggest result. Peak drops from 1.64 GB to 1.13 GB --- 514 MB reduction, 31\% of peak. Only 16 evictions. Three early convolutions are evicted because they produce huge 50 MB feature maps but are relatively cheap. All deeper convolutions retained. 2.6 ms extra cost --- excellent trade-off.}

\qa{Why are early convolutions evicted but deeper ones aren't?}{Early convolutions operate on high-resolution inputs (224x224) with few channels, producing large outputs (50 MB) cheaply (0.23 ms). Deeper convolutions have smaller spatial dimensions but more channels --- the outputs are smaller and the computation is heavier. The ratio favors evicting early layers.}

\slide{Results: BERT-base}
\say{Zero recomputed, 363 retained. AC can't help --- peak is in optimizer. The algorithm correctly detects this.}

\slide{Validation Results}
\say{Three-level validation: profiler accuracy, AC sanity, simulator consistency. All pass on all models. 55 unit tests.}


\section{Wrap-Up (Slides 29--31)} \timing{$\sim$4 min}

\slide{Key Findings}
\say{Five findings. AC works great for vision models. AC doesn't help BERT with foreach Adam. Profiler accuracy scales with tensor size. Memory simulator is essential. Cascading cost propagation matters.}

\slide{Next Steps: Phase 3}
\say{Graph rewriting: extract forward subgraph, clone it, insert before first backward use. The starter code provides utilities. Correctness: torch.allclose on gradients.}

\qa{What's the complexity of graph rewriting?}{Linear in the number of evicted nodes times the size of each recomputation subgraph. For ResNet50 with 16 evictions, this is fast --- the subgraphs are small (1--3 nodes each for relu, transpose, getitem).}

\qa{How does this compare to PyTorch's built-in torch.utils.checkpoint?}{torch.utils.checkpoint requires manual annotation. Our approach is automatic. But it requires FX-traceable models and single-iteration profiling.}

\qa{Could you use DP instead of greedy?}{Yes --- 0-1 knapsack formulation. But greedy is O(N\textsuperscript{2}) vs exponential for exact solutions, and sufficient for our models.}

\slide{Questions}
\say{Happy to take questions or do a live demo.}

\end{document}
